\family{Combretaceae}{Terminalia family}
\label{fam:combretaceae}

Trees, shrubs or climbers. Leaves simple and entire, alternate or
opposite, in \taxon{Terminalia} crowded at the ends of short shoots on
tiered branches. Flowers small, in spikes or heads. Fruit dry and
indehiscent, usually with two to five wings. Two species of
\taxon{Terminalia} in the Arami Uplands, both important timber trees.

\section*{Key to the species}
\begin{florakey}
	\lead{1a}{Bark dark and deeply fissured; buttresses small; fruit
		longer than broad}{\taxon{Terminalia ivorensis}}
	\lead{1b}{Bark pale grey and smooth; tall plank buttresses; fruit much
		broader than long}{\taxon{Terminalia superba}}
\end{florakey}

\species{Terminalia ivorensis}{A.Chev.}
\label{sp:COM01}

\field{Local and trade names} Yoruba: Idigbo. Trade: Idigbo.

\field{Status} Native.

\field{Description} Deciduous tree to 45~m with a straight,
cylindrical bole and small buttresses. Bark dark grey to black, deeply
fissured. Branches in tiers. Leaves crowded at the ends of the shoots,
elliptic to obovate, 5--10~cm long. Flowers small, cream, in spikes.
Fruit with two wings, oblong, longer than broad.

\field{Phenology} Flowering February to April; fruiting April to June.

\field{Habitat} Lowland semi-deciduous forest.

\field{Uses} A light, easily worked timber used for joinery, window
frames and furniture; also planted in forest plantations
\citep{louppe2008}.

\field{Specimens examined} Olowo Stream (D3), 440~m, fl., 3~Apr 1936, \specimen{Harding}{1252}{OKF};
Igbo-Oba Forest Reserve (C3), 470~m, fr., 21~Apr 1937, \specimen{Harding}{1453}{OKF};
Igbo-Oba Forest Reserve (C3), 470~m, fl., 17~Feb 1938, \specimen{Harding}{1563}{OKF};
Igbo-Oba Forest Reserve (C3), 470~m, fl., 3~Apr 2022, \specimen{Alade}{444}{AUH};
Ireke Ferry (D4), 360~m, st., 4~Oct 2023, \specimen{Alade}{574}{AUH}.

\field{Distribution} Grid squares C3, D3, D4.

\begin{center}
	\distmap{C/3, D/3, D/4}
\end{center}

\species{Terminalia superba}{Engl. \& Diels}
\label{sp:COM02}

\field{Local and trade names} Yoruba: Afara. Trade: Afara.

\field{Status} Native.

\field{Description} Deciduous tree to 50~m with tall plank buttresses.
Bark smooth, pale grey. Branches in tiers. Leaves crowded at the ends
of the shoots, obovate, 6--12~cm long. Flowers small, greenish-white,
in spikes. Fruit with two wings, much broader than long.

\field{Phenology} Flowering February to April; fruiting May to August.

\field{Habitat} Lowland semi-deciduous forest and forest--savanna
mosaic, often in secondary forest.

\field{Uses} A pale timber used for plywood, furniture and interior
work, and widely planted in forest plantations
\citep{keay1989,louppe2008}.

\field{Specimens examined} Okeri (C4), 420~m, fr., 12~Jun 1935, \specimen{Harding}{1166}{OKF};
Olowo Stream (D3), 440~m, st., 26~Jan 1937, \specimen{Harding}{1345}{OKF};
Aramiwo Foot (B2), 640~m, fl., 21~Apr 1966, \specimen{Olaniyan}{722}{SVN};
Oke-Asa farmland (B4), 690~m, fr., 4~May 1967, \specimen{Olaniyan}{888}{SVN};
Oke-Asa farmland (B4), 690~m, st., 25~Nov 2021, \specimen{Alade}{394}{AUH}.

\field{Distribution} Grid squares B2, B4, C4, D3.

\begin{center}
	\distmap{B/2, B/4, C/4, D/3}
\end{center}